\documentclass[11pt,a4paper]{article}
\usepackage[margin=25mm]{geometry}
\usepackage[T1]{fontenc}
\usepackage{lmodern,amsmath,booktabs,array}
\usepackage[colorlinks=true,allcolors=blue]{hyperref}
\hypersetup{pdftitle={Submission audit: nearest-neighbour distances and spatial cross-validation},pdfauthor={Prepared for Olatunji Johnson}}
\title{Submission audit}
\author{Nearest-neighbour distances and prediction error\\in spatial cross-validation}
\date{12 November 2026}
\begin{document}
\maketitle
\noindent\textbf{Assessment: the numerical results and central mathematics are intact, but the files are not yet a final submission package.} The main paper and supplementary material compile, the reported numbers agree with the retained results, and the central proofs are coherent. Several presentation and completeness issues have been corrected. The remaining decisions concern the covariance-fitting sensitivity, publication access to code and data, author declarations, and journal formatting.

The draft cover letter names \emph{Environmental and Ecological Statistics}, but calls this choice provisional. Its current submission instructions were used provisionally for the format checks below. No submission or public release has been made.

\section*{Corrections made}
\begin{enumerate}
\item \textbf{Anisotropic neighbour definition.} Section 3.3 and supplementary material Section S5.3.1 now explain that the true-covariance reference uses the most correlated training observation. This is the nearest observation in the known anisotropic metric, and it can differ from the geographically nearest observation used by the isotropic diagnostics. Section S5.3.2 states the corresponding choice for the bound. These changes align the description with the executed code; they do not change results.
\item \textbf{Target locations versus target responses.} Section S5.6.1 now states that target response values are excluded from fitting and selection. Target locations are used by kNNDM and the diagnostics. The earlier wording excluded ``target observations'' from both, which was ambiguous.
\item \textbf{Missing reproduction outputs.} The GitHub snapshot now includes the two fold-count PDFs cited in Section S5.4.1 and the tie-tolerance table cited in Section S5.5. Their generation has been added to the root reproduction workflow. The fold-count report's obsolete cluster-specific instructions were replaced with portable commands. The supplementary material now gives the correct path to the tie table.
\end{enumerate}

\section*{Items requiring action before submission}
\begin{enumerate}
\item \textbf{Assess the covariance-fitting sensitivity.} The 1,440-dataset diagnostic study and 1,920-dataset refitting study use the original variogram fitter. Jura uses a profile search that can find a better solution to the same objective. The estimated-covariance selection gain is small and concentrated in one setting; it has not been checked with that search. This is disclosed in the Discussion and Section S5.6.2. A paired sensitivity rerun would strengthen the empirical claim. Otherwise, the interpretation must remain explicitly limited to the tested fitting procedure. This audit does not establish how much a different fitter would change those results.
\item \textbf{Finalise code and data access.} The paper still says that a permanent public archive remains to be assigned. The GitHub repository is private and the current files are on a draft pull request. Choose reviewer access and the final availability statement; a private link alone does not give reviewers access. A DOI has not been assigned. No repository visibility was changed during this audit.
\item \textbf{Replace internal draft text with final declarations.} The main paper still contains a ``Draft status and assistance'' section. It should become an appropriate factual AI-use disclosure, with author approval handled outside the scientific text. Funding and competing-interest statements have not been supplied; none have been invented. The title pages need a corresponding-author email and complete affiliation. An institutional email is already present in the draft cover letter and can be reused if current.
\item \textbf{Complete journal formatting.} If \emph{Environmental and Ecological Statistics} is retained, its abstract limit is 150--250 words; the current abstract is approximately 270. A 222-word alternative is supplied below for review, without replacing the agreed abstract. The journal requests full DOI links and author/contact details in supplementary files. The current bibliography prints DOI identifiers but not full DOI URLs. Its policy also specifies documentation of substantive LLM use in Methods. These are format checks against the \href{https://link.springer.com/journal/10651/submission-guidelines}{official instructions}, not confirmation of the final journal choice.
\item \textbf{Upload the full supplementary set.} The 18-page supplementary PDF refers to five separate figure PDFs (15 pages in total) and two fold-count PDFs (four pages). They must accompany it, or be consolidated into the final supplementary material. The refreshed reading pack includes all of them. The fold-count report retains author-review branding, which should be removed if it is submitted as a standalone published resource.
\end{enumerate}

\section*{Checks performed}
\begin{itemize}
\item Main paper: 21 pages; supplementary material: 18 pages. Both rebuilt after the wording corrections. No undefined references, unresolved citations or overfull text boxes were found. The remaining underfull-box messages are minor spacing issues.
\item All 22 bibliography entries are cited, with no missing citation keys. DOI metadata and the relevant publisher, repository and data-documentation records were checked. The 2026 target-weighted CV citation is correctly identified as arXiv version 3. The Jura source distinguishes the 359-site distributed dataset from the original survey.
\item The one-neighbour identity, moment decomposition, full-neighbour remainder, and lower/upper bounds were checked algebraically. The supplementary proof includes residual covariances, the quadratic minimisation, the independent-noise eigenvalue bound, the spatial-variance cap, averaging, and boundary cases. No inconsistency was found in those steps.
\item The previous fresh render and simulation checks remain applicable: the numerical code and results were not changed by these clarifications. The updated comparison now checks 37 outputs, including the tie audit; all pass with zero numerical differences. All 132 formatted manuscript values and both generated tables remain unchanged.
\item The primary figures were inspected for legibility, labels and consistency with captions. Figure 1 distinguishes the solid/dashed curves; Figure 4 shows both $W$ and $|J_2|$, as well as the observed target MSE. The main PDF contains all four figures and both tables.
\item Every one of the 37 explicitly named supporting files was located. The two newly restored fold-count PDFs were regenerated, and the 12-row numerical-tolerance table agrees with its retained reference.
\item No bird or PM2.5 application remains in the current main paper or supplementary material. Historical application files are separate from the active reproduction workflow.
\end{itemize}

This is a document, implementation and mathematical consistency audit. It is not a new full-scale simulation rerun or an independent human scientific review. The known covariance-fitting limitation remains unresolved.

\section*{Suggested shorter abstract for author review}
\noindent\textit{This alternative has not replaced the current abstract.}

Spatial cross-validation (CV) can misrepresent prediction error when validation and target locations differ in their relationships with training observations. We propose correlation contrasts to assess these differences, focusing on $J_2$: validation minus target average squared correlation with the nearest training observation. Under Gaussian prediction with known mean and covariance, response variance multiplied by $J_2$ equals target mean squared error (MSE) minus expected CV MSE when prediction uses one neighbour. We derive the additional error reduction supplied by other neighbours and bound its contribution to CV bias. A counterexample shows that matching nearest-distance distributions exactly need not eliminate bias. Among four CV schemes, selection by absolute $J_2$ reduced mean absolute expected bias by 12.9\% relative to selection by Wasserstein distance across 5,120 known-covariance configurations. Across 1,920 datasets with estimated covariance, it reduced the mean absolute error of the observed CV estimate by 2.3\%; almost all of this gain came from one of 24 settings. A ten-fold sensitivity check did not establish a nonzero benefit. Further simulations show that explanatory information can change CV bias while response correlation remains fixed. In the Swiss Jura soil-cadmium application, both selection rules choose kNNDM, but random CV is closer to observed MSE at 100 reserved sites. Correlation contrasts quantify one component of CV bias; their usefulness for scheme selection depends on the prediction model and covariance estimation.
\end{document}
